\documentclass[a4paper,9pt]{extarticle}

%=========================================================
% PACKAGES
%=========================================================

\usepackage[
a4paper,
top=1.8cm,
bottom=1.5cm,
footskip=7mm,
left=1cm,
right=1cm,
headheight=14pt,
headsep=5mm
]{geometry}

\usepackage[T1]{fontenc}
\usepackage[utf8]{inputenc}
\usepackage[french]{babel}

\usepackage[sfdefault]{FiraSans}
\renewcommand{\familydefault}{\sfdefault}

\usepackage{multicol}
\setlength{\columnsep}{0.7cm}

\usepackage{amsmath,amssymb}
\usepackage{titlesec}
\usepackage{enumitem}

\usepackage{tcolorbox}
\tcbuselibrary{skins,breakable}

\usepackage{xcolor}
\usepackage{graphicx}
\usepackage{float}

\usepackage{fancyhdr}
\usepackage[hidelinks]{hyperref}


%=========================================================
% PAGE SETTINGS
%=========================================================

\setlength{\parindent}{0pt}
\setlength{\parskip}{1.5pt}

\setlength{\abovedisplayskip}{4pt}
\setlength{\belowdisplayskip}{4pt}

\setlist[itemize]{
leftmargin=*,
itemsep=1pt,
topsep=1pt
}
\setlist[enumerate]{
leftmargin=*,
itemsep=1pt,
topsep=1pt
}


%=========================================================
% COLORS
%=========================================================

\definecolor{Blue}{HTML}{E3F0FF}
\definecolor{Green}{HTML}{E3F7EA}
\definecolor{Pink}{HTML}{FFE5ED}
\definecolor{Orange}{HTML}{FFF0DC}
\definecolor{Purple}{HTML}{EEE7FF}
\definecolor{Yellow}{HTML}{FFFADB}
\definecolor{Goat}{HTML}{F7EEDF}

\definecolor{BlueTitle}{HTML}{285B8F}
\definecolor{GreenTitle}{HTML}{287052}
\definecolor{PinkTitle}{HTML}{A84063}
\definecolor{OrangeTitle}{HTML}{A96500}
\definecolor{PurpleTitle}{HTML}{6245A5}
\definecolor{GoatTitle}{HTML}{8A5A1F}
\definecolor{RecapTitle}{HTML}{5B3A8E}


%=========================================================
% HEADER
%=========================================================

\pagestyle{fancy}
\fancyhf{}

\fancyhead[L]{\small\textbf{Comportement mécanique des matériaux}}
\fancyhead[C]{\small Weibull, fissuration lente, production de Al et Cu}
\fancyhead[R]{\small Week 06 --- 12-18/10/2026}

\fancyfoot[C]{\small\href{https://github.com/AlphaXT38G/notes_ba5}{github.com/AlphaXT38G/notes\_ba5}}
\fancyfoot[R]{\small \thepage}

\renewcommand{\headrulewidth}{0.4pt}


%=========================================================
% TITLES
%=========================================================

\titleformat{\section}
{\large\bfseries\color{BlueTitle}}
{}{0pt}{}

\titlespacing{\section}
{0pt}{6pt}{3pt}


%=========================================================
% BOXES
%=========================================================
% Boxes are kept whole (unbreakable): inside multicols, a box split by tcolorbox
% leaves an empty column behind it.

\newtcolorbox{definition}{
enhanced, unbreakable, colback=Blue, colframe=BlueTitle,
boxrule=0.7pt, arc=0mm, title=\textbf{Définition}, fonttitle=\bfseries
}

\newtcolorbox{important}{
enhanced, unbreakable, colback=Pink, colframe=PinkTitle,
boxrule=0.7pt, arc=0mm, title=\textbf{Idée clé}, fonttitle=\bfseries
}

\newtcolorbox{formula}{
enhanced, unbreakable, colback=Green, colframe=GreenTitle,
boxrule=0.7pt, arc=0mm, title=\textbf{Formules clés}, fonttitle=\bfseries
}

\newtcolorbox{warning}{
enhanced, unbreakable, colback=Orange, colframe=OrangeTitle,
boxrule=0.7pt, arc=0mm, title=\textbf{Attention}, fonttitle=\bfseries
}

\newtcolorbox{example}{
enhanced, unbreakable, colback=Purple, colframe=PurpleTitle,
boxrule=0.7pt, arc=0mm, title=\textbf{Exemple}, fonttitle=\bfseries
}

% GOAT - brown: the idea told as simply as possible
\newtcolorbox{goat}{
enhanced, unbreakable, colback=Goat, colframe=GoatTitle,
boxrule=0.7pt, arc=0mm, title=\textbf{Expliqué à une chèvre}, fonttitle=\bfseries
}

% RECAP - violet: what to know for the exam, at the end of the week
\newtcolorbox{recap}{
enhanced, breakable, colback=Purple, colframe=RecapTitle,
boxrule=1pt, arc=0mm, title=\textbf{Récap pour l'examen}, fonttitle=\bfseries
}

% Screenshot of a slide of the course: \slide{cNN-PP}
\newcommand{\slide}[2][0.95]{%
\begin{center}
\includegraphics[width=#1\linewidth]{../images/course/#2}
\end{center}}


%=========================================================
% DOCUMENT
%=========================================================

\begin{document}

\thispagestyle{fancy}

\begin{multicols}{2}

%=========================================================
\section{Rupture fragile, résistance statistique}

\begin{important}
Pas de plasticité : rupture au \textbf{plus gros défaut} bien orienté,
$\sigma_r = K_c/(Y\sqrt{\pi a_{\max}})$. Faciès : miroir, zone brumeuse, branches (utile en analyse
d'avarie). Taille des défauts aléatoire $\Rightarrow$ résistance \textbf{statistique}, effet de
\textbf{volume} (gros = moins résistant), effet de la \textbf{répartition} des contraintes.
\end{important}

\slide[0.9]{c02-87}

\begin{goat}
Une chaîne casse à son maillon le plus faible, et une grande chaîne a plus de chances d'en avoir
un mauvais.
\end{goat}

%=========================================================
\section{Statistique de Weibull}

\textbf{D'où elle vient} : volume homogène, un petit volume qui casse fait tout casser ;
$S(V_1+V_2) = S(V_1)S(V_2)$ $\Rightarrow S = e^{-V\Phi(\sigma)}$ ; Weibull pose
$\Phi = K(\sigma-\sigma_u)^m$ (3 paramètres), souvent $\sigma_u = 0$ (2 paramètres).

\begin{formula}
\[ S = 1-F = \exp\!\Big[-\frac{V}{V_0}\Big(\frac{\sigma}{\sigma_0}\Big)^m\Big] \]
\[ \ln S = \ln[1-F] = -\frac{V}{V_0}\Big(\frac{\sigma}{\sigma_0}\Big)^m, \quad \ln[1-F] \approx -F \ (F\ll1) \]
Surface : $A/A_0$. $V_0$, $\sigma_0$ par paire : $K = [V_0\sigma_0^m]^{-1}$.
\end{formula}

\slide[0.8]{c02-95}

\begin{warning}
$\sigma = \sigma_0$, $V = V_0$ : $S = e^{-1}$, \textbf{pas} la médiane.
\end{warning}

\begin{important}
Non uniforme : $S = \exp[-\int_V(\sigma/\sigma_0)^m dV/V_0] = \exp[-k_m\frac{V}{V_0}(\frac{\sigma_{\max}}{\sigma_0})^m]$.
$k_m = 1$ traction ; $1/[2(m+1)]$ flexion pure ; $1/[2(m+1)^2]$ 3 points ;
$(mL_i+L_o)/[2L_o(m+1)^2]$ 4 points. $k_m < 1$ : flexion \textbf{flatteuse}.
\end{important}

\begin{important}
\textbf{Mesure} : classer $\sigma_1 \le \ldots \le \sigma_N$, $F_i \approx i/N$ ;
$\ln\ln(1/S) = m\ln\sigma_{\max} - m\ln\sigma_0 + \ln(k_mV/V_0)$ : droite de pente $m$.
\end{important}

\begin{example}
\textit{(pas dans les diapos)} À même $S$, $V\sigma^m$ constant : pièce 1000$\times$ plus grosse,
$m = 10$ : $\sigma \times 1000^{-0{,}1} = 0{,}5\,\sigma$.
\end{example}

\begin{goat}
Weibull donne la chance qu'aucun maillon ne lâche. $m$ grand : maillons tous pareils, fiable ;
$m$ petit : très inégaux.
\end{goat}

%=========================================================
\section{Compression et fissuration lente}

\begin{formula}
\[ \sigma_{\text{compr.}} = C\frac{K_c}{Y\sqrt{\pi\bar a}}, \quad C \approx 15 \]
Fissures parallèles à la compression, croissance \textbf{stable}, $\bar a$ = longueur
\textbf{moyenne} ; émiettement.
\end{formula}

\begin{formula}
\[ \frac{da}{dt} = F(P_{\text{H}_2\text{O}}, K)\,e^{-Q/RT}, \quad K = Y\sigma\sqrt{\pi a} \]
L'eau casse les Si--O tendues en pointe de fissure ($K < K_c$ : sous-critique, \og fatigue
statique\fg{}) ; seuil à bas $K$, plateau à haut $K$.
\[ \Big(\frac{\sigma}{\sigma_{TS}}\Big)^n = \frac{t_{\text{test}}}{t}, \quad n \approx 10\text{--}20 \]
\end{formula}

\begin{important}
Fibre de verre fraîche : 14\,GPa ; à l'air : $\approx 5$\,GPa.
\end{important}

\begin{goat}
Écrasée, une céramique tient 15 fois mieux. Et l'humidité grignote lentement les fissures
tendues : la vitre finit par casser sous une charge qu'elle tenait au début.
\end{goat}

%=========================================================
\section{Métaux : panorama}

\begin{important}
Un élément principal + éléments d'alliage. Moins rigides et durs que les céramiques, souvent plus
lourds, mais \textbf{ductiles et tenaces} ; oxydables ; empreinte énergie/CO$_2$ notable ;
recyclables. Fe, Al, Ti, Mg, Ni : \textbf{même $E/\rho$}.
\end{important}

\begin{definition}
\textbf{Réserves} : exploitables avec profit aujourd'hui. \textbf{Ressources} : tout ce qu'on
peut envisager d'extraire. Énergie d'un métal : croît avec la dilution et la stabilité de
l'oxyde.
\end{definition}

\begin{formula}
Taux de recyclage : $F \le (1+i)^{-n}$ ($i$ : croissance annuelle, $n$ : durée de vie).
\end{formula}

\begin{goat}
Tant qu'on veut chaque année plus de métal, les vieux objets ne suffisent pas à tout fournir.
\end{goat}

%=========================================================
\section{Production de Cu et Al}

\slide[0.85]{c03a-45}

\begin{definition}
\textbf{Ellingham} : $\Delta G^\circ$ d'oxydation vs $T$. Droite basse = oxyde stable ; un
élément plus bas réduit l'oxyde d'un plus haut ; C + O$_2$ $\to$ 2CO descend avec $T$.
\end{definition}

\begin{important}
\textbf{Cu} : 50\,ppm (rare) mais quasi noble. Flottation, \textbf{matte}, convertisseur (Fe en
scorie, S en SO$_2$), blister, \textbf{raffinage électrolytique}. Aussi voie hydrométallurgique.
\end{important}

\begin{important}
\textbf{Al} : abondant, oxyde très stable. \textbf{Bayer} (bauxite $\to$ Al$_2$O$_3$) puis
\textbf{Hall--Héroult} : électrolyse dans la cryolithe,
$2\text{Al}^{3+} + 6e^- \to 2\text{Al}$. $12\pm2$\,kWh/kg (théorie 6,5) ; recyclage
$\approx 1$\,kWh/kg ; 100\,Mt/an dont 30\,\% recyclé.
\end{important}

\begin{example}
$12$\,kWh/kg $= 43\,200$\,kJ/kg ; $\frac12v^2$ à 100\,km/h $= 386$\,J/kg $\Rightarrow$
112\,000 accélérations, 31 par jour sur 10 ans.
\end{example}

\begin{important}
\textbf{Passivation} : Al$_2$O$_3$ compatible avec le métal (Pilling--Bedworth $\approx 1$),
couche protectrice. Mise en forme : laminage, extrusion, forgeage, fonderie, impression 3D.
\end{important}

\begin{goat}
Ellingham classe qui tient le plus fort à l'oxygène : l'aluminium tellement que seul le courant
l'en sépare ; le cuivre lâche facilement mais il est rare.
\end{goat}

%=========================================================
\section{Récap de la semaine}

\begin{recap}
\begin{itemize}
\item Rupture au plus gros défaut ; effets de volume et de répartition.
\item Weibull : $S = \exp[-(V/V_0)(\sigma/\sigma_0)^m]$, démonstration, $k_m$, flexion flatteuse,
      $\ln\ln(1/S)$ vs $\ln\sigma$, $F \approx$ exposant si $F \ll 1$.
\item Compression $CK_c/(Y\sqrt{\pi\bar a})$, $C \approx 15$.
\item $da/dt = F(P_{\text{H}_2\text{O}},K)e^{-Q/RT}$ ; $(\sigma/\sigma_{TS})^n = t_{\text{test}}/t$.
\item Métaux : même $E/\rho$ ; $F \le (1+i)^{-n}$ ; Ellingham ; Cu (matte, électrolyse) ;
      Al (Bayer, Hall--Héroult, 12\,kWh/kg) ; passivation.
\end{itemize}
\end{recap}

\end{multicols}

\end{document}
